% SGA 3, Expose VI_B, Loop 1 English reconstruction.
% Scope: title and opening through Lemma 1.2.1 in full;
% Counterexample 1.2.2 is excluded.
% Authority: Polo--Gille Exp6B-13oct24.pdf, component-local pages 1--2
% (combined-reader pages 342--343; original-margin printed pages 319--321).
% OCR and the Reinhold English Markdown are comparison witnesses only.

\setcounter{footnote}{-1}
\chapter*{第 \texorpdfstring{VI\textsubscript{B}}{VI_B} 篇\\群概形的一般性质}
\phantomsection
\label{sga3:VIB:title}
\addcontentsline{toc}{chapter}{第 VI-B 篇. 群概形的一般性质}
\markboth{第 VI\textsubscript{B} 篇. 群概形的一般性质}%
  {第 VI\textsubscript{B} 篇. 群概形的一般性质}

\begin{center}
  \textit{J.-E. Bertin 著}
\end{center}

\footnote{\label{note:sga3-VIB-0}编者注: 2024 年 10 月 13 日版本.}
本篇不对应于讨论班中的任何一次口头报告,
目的是汇集关于群概形的一些常用技术性结果.%
\begingroup
\renewcommand{\thefootnote}{\fnsymbol{footnote}}%
\footnote[1]{本篇在油印版出版之后作了相当大的修订;
特别地, \hyperref[sec:sga3-VIB-10]{第 10 节}和第 11 节完全重写了.
\textsuperscript{(1)}}%
\endgroup
\footnotetext[1]{编者注: 还有 \hyperref[sec:sga3-VIB-5]{第 5 节};
参见编者\hyperref[note:sga3-VIB-34]{注 34} 和注 46.}
\setcounter{footnote}{1}

\section*{1. 域上局部有限型群的态射}
\phantomsection
\label{sec:sga3-VIB-1}
\addcontentsline{toc}{section}{1. 域上局部有限型群的态射}

\medskip
\phantomsection\label{par:sga3-VIB-1.1}%
\noindent\textbf{1.1.---}
设 \(A\) 为 Artin 局部环, \(G\) 与 \(H\) 为两个 \(A\)-群,%
\footnote{\label{note:sga3-VIB-2}编者注: 回顾第~VI\(_A\)~篇开头
采用的约定: \(A\)-群是 \(A\) 上的群概形
(且每个这样的概形都是分离的; 参见第~VI\(_A\)~篇,
\SGARefLink{sga3:VIA:0.3}{0.3}).}
\(u:G\to H\) 为 \(A\)-群态射. 则 \(u\) 诱导群态射
\[
  u(A):G(A)\longrightarrow H(A).
\]%
\footnote{\label{note:sga3-VIB-3}编者注: 已将``同态''替换为``态射''.}
由于 \(H(A)\) 通过右平移作用在 \(H\) 上,
沿 \(u(A)\) 限制就定义了 \(G(A)\) 在 \(H\) 上的作用.
这一作用与 \(u\) 以及 \(G(A)\) 在 \(G\) 上的右平移作用相容.
由于 \(G(A)\) 在 \(G\) 的严格有理点上传递地作用
(定义参见第~VI\(_A\)~篇, \SGARefLink{sga3:VIA:0.4}{0.4}),
这些点``关于 \(u\) 都有相同的表现''.
以下性质由这一观察产生.

\medskip
\phantomsection\label{prop:sga3-VIB-1.2}%
\noindent\textbf{命题 1.2.---}
\emph{设 \(u:G\to H\) 是 \(A\) 上局部有限型的
\(A\)-群之间的拟紧态射. 则集合 \(u(G)\) 在 \(H\) 中闭, 且}
\[
  \dim G=\dim u(G)+\dim\ker u.
\]%
\footnote{\label{note:sga3-VIB-4}编者注: 这一陈述也见于第~VI\(_A\)~篇,
\SGARefLink{prop:sga3-VIA-2.5.2}{2.5.2}.}

由于 \(u\) 与 \(G\) 和 \(H\) 的取逆态射可交换,
像 \(u(G)\) 在 \(H\) 的取逆作用下不变;
因此它在 \(H\) 中的闭包 \(\overline{u(G)}\) 也如此.
另一方面, 设 \(L\) 为 \(H\times_A H\) 中两个投影均属于 \(u(G)\)
的点所组成的集合. 显然, \(L\) 是
\[
  u\times_Au:G\times_A G\longrightarrow H\times_A H.
\]
的像. 因而 \(H\) 的乘法态射将 \(L\) 映入 \(u(G)\); 换言之,
\[
  u(G)\cdot u(G)=u(G).
\]
下面的 \hyperref[lem:sga3-VIB-1.2.1]{引理~1.2.1} 表明,
\(L\) 在 \(H\times_AH\) 中的闭包 \(\overline L\),
是两个投影均属于 \(\overline{u(G)}\) 的点所组成的集合. 因而
\[
  \overline{u(G)}\cdot\overline{u(G)}
    =\overline{u(G)}.
\]
于是, \(H\) 中以 \(\overline{u(G)}\) 为底空间的约化子概形,
自然带有范畴 \((\mathbf{Sch}_{/k})_{\mathrm{red}}\) 中的群结构,
其中 \(k\) 是 \(A\) 的剩余域
(参见第~VI\(_A\)~篇, \SGARefLink{sga3:VIA:0.2}{0.2}).

证明 \hyperref[prop:sga3-VIB-1.2]{1.2} 的第一个断言.
将 \(A\) 替换为其剩余域 \(k\) 的一个代数闭包后,
可假设 \(A\) 是代数闭域 \(k\) (参见 EGA~IV\(_2\), 2.3.12).
将 \(u\) 替换为
\[
  u_{\mathrm{red}}:G_{\mathrm{red}}\longrightarrow H_{\mathrm{red}},
\]
可假设 \(G\) 和 \(H\) 约化. 如刚才所见,
此时 \(u(G)\) 是 \(H\) 的一个约化群子概形的底空间;%
\begingroup
\renewcommand{\thefootnote}{\fnsymbol{footnote}}%
\footnote[1]{\label{note:sga3-VIB-1.2-image-ambient}译者注:
法文底本在此印作``\(G\) 的一个约化群子概形''.
由于 \(u(G)\subset H\), 且前一段是在 \(H\) 内构造约化子概形,
这里的环境群应为 \(H\).}%
\endgroup{}
因而可假设 \(u\) 为支配态射. 此时 \(G(k)\) 在 \(H\)
的连通分支集上传递地作用, 只须证明 \(u(G)\cap H^0\) 闭.
于是化归到 \(H\) 连通的情形, 从而它不可约且为有限型
(第~VI\(_A\)~篇, \SGARefLink{cor:sga3-VIB-2.4}{2.4}).
此时态射 \(u\) 为有限型, 因为它拟紧且局部有限型.
由于 \(H\) 是 Noether 概形, \(u(G)\) 可构造
(EGA~IV\(_1\), 1.8.5), 因而包含 \(H\) 的一个开子集 \(V\)
(EGA~0\(_{\mathrm{III}}\), 9.2.2).
由第~VI\(_A\)~篇, 0.5, 有
\[
  H=V\cdot V\subset u(G)\cdot u(G)=u(G).
\]

证明第二个断言. 先回顾一下, 函子 \(\ker u\)
(参见 \SGARefLink{sga3:i:definition:2-3-6-1}{第~I~篇, 2.3.6.1})
由 \(u^{-1}(e)\) 表示, 其中 \(e\) 表示 \(H\) 的单位元.
因此当 \(u\) 局部有限型时, \(\ker u\) 在 \(A\) 上局部有限型.
与上面相同, 这次使用 EGA~IV\(_2\), 4.1.4,
将问题化归到 \(A\) 为代数闭域 \(k\) 的情形.

还可假设 \(G\) 和 \(H\) 不可约且为有限型, \(u\) 为支配态射.
事实上, 由于 \(k\) 代数闭, \(G\) 的连通分支显然都有相同维数,
因为 \(G(k)\) 在这些分支构成的集合上传递地作用;
若用 \(u^0\) 表示 \(u\) 在 \(G^0\) 上的限制, 则
\[
  (\ker u)^0\subset\ker u^0,
  \qquad
  \dim\ker u^0=\dim\ker u.
\]
同样地, \(u\) 此时在 \(k\) 上为有限型.
若 \(\eta\) 为 \(H\) 的泛点, 则
\[
  \dim u^{-1}(\eta)=\dim G-\dim H
\]
(EGA~IV\(_3\), 10.6.1(ii)). 根据 EGA~IV\(_3\), 9.2.3 和 9.2.6,
满足
\[
  \dim u^{-1}(y)=\dim u^{-1}(\eta)
\]
的 \(y\in H\) 组成的集合包含一个非空开子集 \(V\).
由于 \(u\) 支配, \(U=u^{-1}(V)\) 为 \(G\) 的非空开子集;
它包含一个闭点 \(x\), 因为 \(G\) 是 Jacobson 概形
(EGA~IV\(_3\), 10.4.7). 右平移 \(r_x\) 是从 \(\ker u\)
到 \(u^{-1}(u(x))\) 的同构, 因而
\[
  \dim\ker u
    =\dim u^{-1}(u(x))
    =\dim u^{-1}(\eta)
    =\dim G-\dim H.
\]

\medskip
\phantomsection\label{lem:sga3-VIB-1.2.1}%
\noindent\textbf{引理 1.2.1.---}
\emph{设 \(f:X'\to X\) 与 \(g:Y'\to Y\) 是 Artin 局部环 \(A\)
上的概形之间的两个拟紧支配态射. 则}
\[
  f\times_Ag:X'\times_A Y'\longrightarrow X\times_A Y
\]
\emph{支配 (并且拟紧).}

事实上,
\[
  f\times_Ag
    =
  (\operatorname{id}_X\times_Ag)
    \circ
  (f\times_A\operatorname{id}_{Y'}).
\]
因此只须证明 \(f\times_A\operatorname{id}_{Y'}\) 与
\(\operatorname{id}_X\times_Ag\) 支配.
为此, 可将 \(A\) 替换为其剩余域 \(k\).
于是 \(X\) 和 \(Y'\) 在 \(A=k\) 上平坦,
而 \(f\times_A\operatorname{id}_{Y'}\)
(分别地, \(\operatorname{id}_X\times_Ag\))
由 \(f\) (分别地, \(g\)) 经平坦基变换 \(Y'\to A\)
(分别地, \(X\to A\)) 得到.
故由 EGA~IV\(_2\), 2.3.7, 它支配 (且拟紧).
% Translation statistics: source_lines=205; translated_lines=187; inline_math=134; display_math=13; pending_TODO=0.
% Review: complete component translation; math fragments, reference keys, environments and note counts checked against source; formula prose translated; compilation pending parent integration.
